% T32 source: Примеры базовой модели.tex lines 822--891
\detailedexampleheading{Comparison of Variants}

\begin{table}[ht!]
\centering
\caption{Comparison of work-organization variants accounting for the dependency graph}
\label{tab:practical-comparison}
\small
\begin{tabularx}{\textwidth}{Xcccc}
\toprule
Metric & Variant 1 & Variant 2 & Variant 3 & Variant 4 \\
\midrule
Configured parallelism $P$ & 1 & 2 & 4 & 4 \\
Mode & interactive & mixed & delegated w/spec. & delegated w/spec. \\
Decomposition & coarse & coarse & coarse & split tests \\
$W$, h & 65.0 & 59.8 & 53.6 & 53.6 \\
$W/P$, h & 65.0 & 29.9 & 13.4 & 13.4 \\
$L$, h & 51.8 & 47.8 & 43.2 & 36.7 \\
$H$, h & 65.0 & 38.4 & 19.0 & 19.0 \\
$B_4$, h & 65.0 & 47.8 & 43.2 & 36.7 \\
$T(\pi)$, h & 65.0 & 47.8 & 43.2 & 36.7 \\
$S_E^{\text{ach}}$ & $\times 1.17$ & $\times 1.59$ & $\times 1.76$ & $\times 2.07$ \\
Ceiling $1/\bar{h}_w$ & $\times 1.17$ & $\times 1.98$ & $\times 4.00$ & $\times 4.00$ \\
\bottomrule
\end{tabularx}
\end{table}

For these illustrative variants, the table shows that, once dependencies are
accounted for, the final speedup is determined not only by the total work $W$
and the number of agents $P$, but also by the critical-path length $L$---and
$L$ itself depends on the selected decomposition.

\begin{enumerate}
  \item \textbf{The $W/P$ lower bound is informative but optimistic as a
  makespan estimate.} It equals $29.9$~h for Variant~2 and $13.4$~h for
  Variant~3. These values are attainable only with sufficiently independent
  tasks and appropriately timed requests for human attention; queueing
  increases agent-stream occupancy beyond $W$.

  \item \textbf{The critical path determines the result in Variants 2--4.} In
  Variants 2 and 3, the chain $1 \to 2 \to 3 \to 4 \to 6$ determines
  makespans of $47.8$ and $43.2$~h, respectively. After the tests are split in
  Variant 4, the critical path decreases to $36.7$~h but remains the active
  constraint. The achievable speedups are therefore $\times 1.59$,
  $\times 1.76$, and $\times 2.07$, rather than the values calculated under
  the assumption of fully independent tasks.

  \item \textbf{Increasing $P$ alone is insufficient in these illustrative
  variants.} In the transition from $P=2$ to $P=4$, the independent-task
  estimate $W/P$ decreases sharply, but the overall makespan changes only
  moderately: additional agents cannot begin tasks on the critical path until
  their predecessors have been completed.

  \item \textbf{The decomposition contrast is isolated by an all-else-equal
  comparison.} The transitions $1 \to 2$ and $2 \to 3$ change several
  configuration parameters simultaneously. Only the transition $3 \to 4$
  preserves $P$, the operating mode, $k$, and $h$; therefore, within this
  illustrative scenario, the increase in speedup from $\times 1.76$ to
  $\times 2.07$ pertains to the change in decomposition and requires no
  additional resources.

  \item \textbf{The human resource imposes a ceiling and constrains the
  schedule.} As the human component decreases, the ceiling $1/\bar{h}_w$
  increases: $\times 1.17 \to \times 1.98 \to \times 4.00$. It is already
  attained in Variant 1, where the process is human-constrained. Variants 2--4
  retain headroom, but exploiting it requires explicit scheduling of the
  service intervals. In the revised Variant 2 schedule, the agent assigned to
  Task 7 is blocked in the queue for developer attention for $22.2$~h; this
  does not increase the makespan because the task is not on the critical path.

  \item \textbf{Moderate attention-switching overhead does not change the
  ranking in this example.} With $\gamma(P)=1+0.1(P-1)$, exact recomputation on
  the refined time grid yields the following values (hours):
  \[
  \begin{array}{c|ccccc}
  \text{variant} & W_4 & L_4 & \gamma H & B_4 & T^* \\ \hline
  2 & 63.640 & 51.320 & 42.240 & 51.320 & 51.36 \\
  3 & 59.300 & 47.700 & 24.700 & 47.700 & 47.70 \\
  4 & 59.300 & 40.450 & 24.700 & 40.450 & 40.45
  \end{array}
  \]
  Thus, the ordering $4<3<2$ is preserved, and the human-resource branch
  remains inactive. Under the more conservative assumption $h_i=0.5k_i$ for
  Variants 3--4, we obtain $H=26.8$~h and $\gamma(4)H\approx34.8$~h; for
  Variant 4, this value is already close to the critical-path length.
  Consequently, as $L_4$ decreases, the human resource may become the active
  constraint.
\end{enumerate}

\paragraph{Implications of the Illustrative Scenario.}
\begin{enumerate}
  \item \textbf{Dependencies should be modeled explicitly.} Balancing
  independent tasks overstates speedup when the work is linked by a dependency
  graph. For the illustrative scenario considered here, at level M4 the lower
  bound is $\max\{W/P, L, H\}$ when $C(P) = \gamma(P) = 1$.

  \item \textbf{The critical path becomes the principal constraint.} In
  Variants~2--4, the makespan equals $L$. The composition of the critical path
  depends on the granularity of the graph: after the testing task is split,
  only retry testing remains at the end of the path, whereas unit and
  integration tests are performed earlier and in parallel.

  \item \textbf{In this illustrative scenario, increasing the number of agents
  does not yield proportional speedup.} The mechanism illustrated here is that,
  if the graph contains large sequentially dependent tasks, increasing $P$
  quickly ceases to affect the result: some agents remain idle while waiting
  for mandatory predecessors to be completed.

  \item \textbf{Decomposition produces a quantifiable difference in this
  illustrative scenario.} Splitting one testing task reduced the makespan by
  $6.5$~h and increased speedup from $\times 1.76$ to $\times 2.07$ without
  changing $P$, $k$, or $h$. Further candidates should be sought among
  critical-path edges that connect entire large tasks even though the actual
  dependency pertains to only one subtask.

  \item \textbf{Both agents and human involvement must be scheduled.} A single
  developer reviews delegated tasks sequentially, so the corresponding
  intervals and the queue for developer attention must be included explicitly
  in the M4 schedule. A waiting agent retains its assigned agent slot; that
  slot cannot be reassigned to another task until the current task is accepted.
  Prolonged continuous interactive work blocks the processing of parallel
  results, and the achievable speedup is bounded above by $1/\bar{h}_w$
  independently of the number of agents.
\end{enumerate}

\ifdefined\FULLARTICLE
  \let\finishdocument\relax
\else
  \def\finishdocument{\end{document}}
\fi
\finishdocument
